% Part: history
% Chapter: biographies
% Section: emmy-noether

\documentclass[../../../include/open-logic-section]{subfiles}

\begin{document}

\olfileid{his}{bio}{noe}

\olsection{إيمي نويثر}

\olphoto{noether-emmy}{Emmy Noether}

وُلدت إيمي نويثر (\textsc{ner}-ter) في إرلنغن بألمانيا في 23 مارس 1882،
في أسرة علمية من الطبقة المتوسطة العليا. وقد لُقبت «أم الجبر الحديث»،
وقدمت إسهامات رائدة في الرياضيات والفيزياء معاً، على الرغم من العقبات
الكبيرة التي اعترضت تعليم النساء. ففي ألمانيا آنذاك، كان يُفترض أن تتعلم
الفتيات الفنون، ولم يكن يُسمح لهن بالالتحاق بالمدارس الإعدادية للجامعة.
غير أن نويثر، بعد حضورها مقررات بصفة مستمعة في جامعتي غوتنغن وإرلنغن
(حيث كان والدها أستاذاً للرياضيات)، استطاعت في نهاية المطاف التسجيل طالبةً
في إرلنغن عام 1904، حين عُدلت سياسة الجامعة لتسمح بقبول الطالبات. ونالت
الدكتوراه في الرياضيات عام 1907.

واجهت نويثر، على الرغم من مؤهلاتها، مقاومةً شديدة خلال مسيرتها. فدرّست في
إرلنغن من 1908 إلى 1915 من دون أجر. وخلال هذه الفترة استرعت انتباه ديفيد
هيلبرت، أحد أبرز رياضيي العالم آنذاك، فدعاها إلى غوتنغن. لكن النساء كن
ممنوعات من نيل الأستاذية، فلم تستطع المحاضرة إلا باسم هيلبرت، ومن دون أجر
مرة أخرى. وفي تلك الفترة برهنت ما يُعرف اليوم بمبرهنة نويثر، التي لا تزال
مستخدمةً في الفيزياء النظرية. وأُجيز لها التدريس أخيراً عام 1919. ويُروى
أن رد هيلبرت على استمرار مقاومة زملائه في الجامعة كان: «أيها السادة، إن
مجلس الكلية ليس حمّاماً عاماً».

ركزت في أواخر عشرينيات القرن العشرين على الجبر المجرد، وأحدثت إسهاماتها
ثورةً في هذا المجال. وكانت كثيراً ما تستخدم في براهينها ما يسمى شرط
السلسلة الصاعدة، الذي ينص على عدم وجود سلسلة لانهائية متزايدة بصرامة من
مجموعات معينة. فمثلاً، تتصف بعض البنى الجبرية المعروفة الآن بالحلقات
النويثرية بأنه لا توجد فيها متتاليات لانهائية من المثاليات
$I_1 \subsetneq I_2 \subsetneq \dots$. ويمكن تعميم الشرط على أي ترتيب
جزئي (ويتعلق في الجبر بالحالة الخاصة للمثاليات المرتبة بعلاقة الاحتواء)،
كما يمكن النظر في شرط السلسلة النازلة الثنائي، حيث تنتهي في نهاية المطاف
كل متتالية \emph{متناقصة} بصرامة في ترتيب جزئي. وإذا حقق ترتيب جزئي شرط
السلسلة النازلة، أمكن استعمال الاستقراء على امتداد هذا الترتيب بطريقة
تشبه استعمالنا الاستقراء على امتداد الترتيب $<$ على~$\Nat$. وتسمى هذه
الترتيبات \emph{حسنة التأسيس} أو \emph{نويثرية}، ويسمى مبدأ البرهان
المناظر \emph{الاستقراء النويثري}.

كانت نويثر يهودية، وحين تولى النازيون السلطة عام 1933 فُصلت من منصبها.
ولحسن الحظ، استطاعت الهجرة إلى الولايات المتحدة لشغل منصب مؤقت في كلية
برين ماور بولاية بنسلفانيا. وخلال إقامتها هناك حاضرت أيضاً في برينستون،
مع أنها وجدت الجامعة غير مرحبة بالنساء \citep[81]{Dick1981}. وفي عام
1935 خضعت نويثر لعملية استئصال ورم رحمي. وتوفيت من عدوى نتجت عن الجراحة،
ودُفنت في برين ماور.

\begin{reading}
لسيرة نويثر، انظر \citet{Dick1981}. ولدى معهد بيريميتر للفيزياء النظرية
محاضراته عن حياة نويثر وتأثيرها متاحةً على الإنترنت
\citep{Perimeter2015}. وإذا سئمت القراءة، فلدى \emph{أشياء فاتتك في حصة
التاريخ} تدوينة صوتية عن حياتها وتأثيرها \citep{Frey2015}. وأعمال نويثر
المجموعة متاحة بالألمانية الأصلية \citep{Noether1983}.
\end{reading}

\end{document}
